\phantomsection
\section*{ВВЕДЕНИЕ}
\addcontentsline{toc}{section}{Введение}

Подшипники скольжения являются одними из наиболее распространённых узлов трения в современном машиностроении. Они широко применяются в турбинах, компрессорах, насосах и других агрегатах нефтегазового комплекса, где обеспечивают опору вращающихся валов при высоких нагрузках и скоростях. Работоспособность подшипника скольжения определяется условиями формирования смазочного слоя между валом и втулкой: при достаточной скорости вращения и вязкости смазки возникает гидродинамическое давление, разделяющее поверхности и предотвращающее их непосредственный контакт.

Надёжность и долговечность подшипниковых узлов непосредственно влияют на работоспособность всего оборудования. Выход из строя подшипника в нефтегазовом оборудовании может привести к аварийной остановке, значительным экономическим потерям и экологическим последствиям. В связи с этим задача повышения эксплуатационных характеристик подшипников скольжения -- увеличения несущей способности, снижения потерь на трение и улучшения условий смазки -- является актуальной научно-технической проблемой.

Одним из перспективных направлений повышения эксплуатационных характеристик подшипников скольжения является нанесение на рабочую поверхность дискретного микрорельефа -- системы регулярных углублений микронного масштаба. Технология лазерного текстурирования поверхностей, получившая развитие в последние десятилетия, позволяет с высокой точностью формировать углубления заданной геометрии на рабочих поверхностях трибологических узлов. Исследования показали, что правильно подобранный микрорельеф способен снизить коэффициент трения на 10-30\%, увеличить несущую способность смазочного слоя и улучшить теплоотвод.

Физический механизм действия микрорельефа основан на создании локальных клиновых зазоров, генерирующих дополнительное гидродинамическое давление (микрогидродинамический эффект). Каждое углубление на рабочей поверхности представляет собой миниатюрный элемент, на входной кромке которого зазор сужается, создавая локальный рост давления. Кроме того, углубления выполняют функцию микрорезервуаров смазки, обеспечивающих подачу масла при пуске и останове, а также служат ловушками для продуктов износа, снижая абразивное воздействие на рабочие поверхности.

Для количественного анализа влияния микрорельефа на характеристики подшипника применяется численное решение уравнения Рейнольдса -- основного уравнения гидродинамической теории смазки. Уравнение Рейнольдса описывает распределение давления в тонком слое вязкой жидкости и позволяет рассчитать интегральные характеристики подшипника: нагрузочную способность, коэффициент трения и расход смазки. Численное решение методом конечных разностей с итерационной процедурой Гаусса-Зейделя обеспечивает получение поля давления на расчётной сетке с учётом граничных условий и кавитации.

Актуальность работы обусловлена потребностью нефтегазовой отрасли в повышении надёжности и энергоэффективности оборудования. Применение текстурированных поверхностей в подшипниках скольжения позволяет улучшить трибологические характеристики без существенного изменения конструкции узла, что делает данный подход привлекательным с практической точки зрения.

Цель курсовой работы -- освоение методов анализа и расчёта трибологических характеристик радиального подшипника скольжения с дискретным микрорельефом.

Задачи:
\begin{enumerate}
    \item Изучить теоретические основы гидродинамической смазки и влияния микрорельефа на характеристики подшипников скольжения.
    \item Освоить методику численного решения уравнения Рейнольдса методом конечных разностей с итерационной процедурой Гаусса-Зейделя.
    \item Выполнить расчёт полей давления и интегральных характеристик для гладкого и текстурированного подшипников в диапазоне эксцентриситетов.
    \item Провести сравнительный анализ результатов и оценить эффективность текстурирования.
\end{enumerate}

Объект исследования -- радиальный подшипник скольжения с дискретным микрорельефом рабочей поверхности.

Предмет исследования -- влияние геометрических параметров микроуглублений на нагрузочную способность, коэффициент трения и расход смазки подшипника.

Методы исследования: численное моделирование на основе метода конечных разностей, итерационный метод Гаусса-Зейделя с верхней релаксацией, сравнительный анализ результатов расчёта.

Работа состоит из введения, трёх глав, заключения, списка использованных источников и приложений. В первой главе представлен обзор литературы по гидродинамической теории смазки и технологии текстурирования поверхностей. Во второй главе описана математическая модель и численный метод решения. В третьей главе приведены результаты расчёта и сравнительный анализ характеристик гладкого и текстурированного подшипников.
